\documentclass[11pt]{article}
\usepackage[T1]{fontenc}
\usepackage[utf8]{inputenc}
\usepackage{lmodern}
\usepackage[a4paper,margin=26mm,headheight=14pt]{geometry}
\usepackage{amsmath,amssymb,amsthm,mathtools,mathrsfs}
\usepackage{microtype,booktabs,longtable,array,enumitem}
\usepackage[dvipsnames]{xcolor}
\usepackage{hyperref,bookmark,fancyhdr}
\definecolor{ResearchBlue}{HTML}{173D63}
\definecolor{ResearchTeal}{HTML}{126B72}
\hypersetup{colorlinks=true,linkcolor=ResearchBlue,citecolor=ResearchTeal,urlcolor=ResearchTeal,pdftitle={Cubic Tornheim Derivatives and Harmonic Stieltjes Constants},pdfauthor={Research continuation for the ProveIt project}}
\pagestyle{fancy}
\fancyhf{}
\fancyhead[L]{\small\itshape Cubic Tornheim and harmonic Stieltjes identities}
\fancyhead[R]{\small ProveIt research}
\fancyfoot[C]{\thepage}
\renewcommand{\headrulewidth}{0.3pt}
\setlength{\emergencystretch}{3em}
\setlength{\parskip}{3pt}
\setlist{topsep=4pt,itemsep=3pt,parsep=0pt}
\allowdisplaybreaks[1]
\numberwithin{equation}{section}
\newtheorem{theorem}{Theorem}[section]
\newtheorem{proposition}[theorem]{Proposition}
\newtheorem{lemma}[theorem]{Lemma}
\newtheorem{corollary}[theorem]{Corollary}
\theoremstyle{definition}
\newtheorem{definition}[theorem]{Definition}
\newtheorem{example}[theorem]{Example}
\theoremstyle{remark}
\newtheorem{remark}[theorem]{Remark}
\newcommand{\Li}{\operatorname{Li}}
\newcommand{\FP}{\operatorname{FP}}
\newcommand{\CT}{\operatorname{CT}}
\newcommand{\PP}{\operatorname{PP}}
\newcommand{\Res}{\operatorname{Res}}
\newcommand{\T}{\mathcal T}
\newcommand{\Z}{\mathbb Z}
\newcommand{\N}{\mathbb N}
\newcommand{\C}{\mathbb C}
\newcommand{\Q}{\mathbb Q}
\newcommand{\Repart}{\operatorname{Re}}
\newcommand{\ii}{\mathrm i}
\newcommand{\dd}{\,\mathrm d}
\newcommand{\etah}{\eta}
\newcommand{\gH}{\gamma_H}
\newcommand{\ray}{W}
\newcommand{\sourcecommit}{\texttt{5a790187c8e186e41e2b990b4941cb7a1a3c7b6b}}

\title{\color{ResearchBlue}\bfseries Cubic Tornheim Derivatives and\\Harmonic Stieltjes Constants\\[6pt]\large A two-coordinate formula, a digamma-cube bridge,\\and finite reductions at mixed integer orders}
\author{Research continuation for the ProveIt project}
\date{11 October 2026 (UTC)}
\begin{document}
\maketitle
\thispagestyle{empty}
\begin{abstract}
This report develops three connected exact-identity results. First, a
coordinate-normalized integral of the digamma cube is expressed through the
first strict harmonic Stieltjes coefficient, linking an ordered Hurwitz
Laurent coordinate to the third derivative of diagonal Tornheim zeta.
A normally convergent harmonic series and an explicitly normalized
antiderivative supply independent analytic descriptions. Second, every
admissible third directional derivative of Tornheim zeta at its singular
origin is reduced to two fixed constants and ordinary Riemann-zeta
derivatives. The coefficient of each nonclassical coordinate is explicit;
an asymmetric conic of directions eliminates one coordinate. A transverse
curved version retains these two coordinates and has an exact fourth-jet
correction. Third, products containing arbitrary positive and nonpositive
integer polylogarithm orders admit finite nested-polylogarithm reductions
with a free complex outer order. The construction includes integer
weights, repeated colors, all outer-order derivatives, and the logarithmic
Mellin poles that distinguish the mixed sector from the purely rational
one. Proofs, exact symbolic certificates, independent numerical checks,
source-status notes, and further research questions accompany the results.
No arithmetic independence of the displayed constants is asserted.
\end{abstract}
\vfill
\noindent\textbf{Source snapshot.} The manuscript and incoming reports were
checked against commit\\[3pt]
\sourcecommit.\\[5pt]
The accompanying source inventory records the Git blob hashes of all
80 manuscript chapter files and all 11 incoming archives at that revision.
\clearpage
\tableofcontents
\clearpage
\section{Research scope and conventions}
\label{sec:scope}

The current ProveIt manuscript already contains extensive exact
polylogarithm, Gamma, Stieltjes, and correlation identities. Its newest
incoming reports go further: they construct arbitrary-depth ordered
Hurwitz germs, complete coincident Stieltjes products, and compute the
first two derivatives of Tornheim zeta on every admissible ray through
the origin \cite{Repo,IncomingOrdered,IncomingCoincident}. The questions
addressed here are specific continuations of that work.

\subsection{What is proved here}

\begin{enumerate}
\item The cubic digamma finite part is reduced to the strict harmonic
coefficient $\eta(1)$ already used in the ordered Hurwitz report:
\[
 \FP\int_0^1\psi(x)^3\dd x
       =3\zeta(2)-6\eta(1)-6\gamma\gamma_1.
\]
This identifies two previously separate coordinates in the inspected
reports. It supplies a convergent harmonic-series answer to the proposed
reduction of the diagonal Tornheim third derivative. It does not give a
finite evaluation of $\eta(1)$ using only ordinary zeta constants.
\item The entire cubic directional layer of $\T(A,B,C)$ is determined
by two fixed constants. An exact double-zeta product identity reduces
the number of unknown quadratic coefficients in its holomorphic
remainder; the resulting formula answers the asymmetric-ray question
in Section~9, Question~2 of the coincident report.
\item The mixed integer-order problem in Section~9, Question~11 of
that report has a finite constructive solution at arbitrary depth.
The free outer spectral order is retained throughout. The pole
description includes logarithmic singularities, which are absent
from the purely rational case.
\end{enumerate}

The Gaussian $S_6$ and revised $S_8$ candidates are different arithmetic
questions. They are not proved by the results below. The exact-identity
and source-status ledgers retain this distinction.

\subsection{Classical input and the extent of novelty}

The Mellin representation, Jonqui\`ere expansion, shuffle and stuffle
products, Mittag--Leffler methods, and the elementary relationship
between Hurwitz zeta and polygamma functions are classical
\cite{DLMFpolylog,DLMFhurwitz,DLMFgamma,Connon,BorweinDilcher}.
General Laurent algorithms for multiple zeta functions also predate
this report \cite{MOW}. Harmonic Stieltjes coefficients have an
established normalization and literature \cite{Kargin}.

The contribution is the specific cubic bridge, the complete formula
for arbitrary cubic directions and transverse curves, and the explicit
finite reduction and pole prescription for the mixed integer-order
sector. These are proved directly. They are new continuations relative
to the inspected repository corpus; a claim of global priority for
every equivalent formulation would require a broader literature audit.
In particular, a two-coordinate spanning formula is not a proof that
the corresponding numbers are arithmetically independent.

\subsection{Notation}

The generalized Stieltjes convention is
\begin{equation}\label{eq:stieltjes}
 \zeta(1+u,a)=\frac1u+
       \sum_{m=0}^{\infty}\frac{(-1)^m\gamma_m(a)}{m!}u^m,
 \qquad \gamma_m=\gamma_m(1),\quad \gamma=\gamma_0.
\end{equation}
Thus $\gamma_0(a)=-\psi(a)$ and $\psi=\Gamma'/\Gamma$.
Primes on $\zeta(s,a)$ denote derivatives in its first argument unless
explicitly stated otherwise. We set
\[
 L=\log(2\pi),\qquad
 H_n^{(r)}=\sum_{j=1}^n j^{-r},\qquad H_n=H_n^{(1)},\qquad H_0^{(r)}=0.
\]
Powers of positive real numbers use real logarithms. All logarithms
of complex variables in a right half-plane use the branch real on
the positive axis.

Our ordered double zeta convention places the larger summation index
first:
\begin{equation}\label{eq:Z2def}
 Z_2(s,t;a)=\sum_{n>m\ge0}(n+a)^{-s}(m+a)^{-t},\qquad a>0.
\end{equation}
The series converges, for example, if $\Re s>1$ and
$\Re(s+t)>2$. Its meromorphic continuation is understood when a
spectral argument is outside this domain. The strict harmonic coefficient
is fixed by
\begin{equation}\label{eq:eta-definition}
 Z_2(1+u,1;a)=\frac1{u^2}+\frac{\gamma_0(a)}u+
 \frac{\gamma_0(a)^2-\zeta(2,a)}2+\eta(a)u+O(u^2).
\end{equation}
This is the convention of the ordered incoming report. If
$\zeta_H(s,a)$ uses the inclusive inner sum, then
\begin{equation}\label{eq:strict-inclusive}
 \zeta_H(s,a)=Z_2(s,1;a)+\zeta(s+1,a),\qquad
 \eta(a)=-\gamma_H(1,a)-\zeta'(2,a).
\end{equation}
The correction $\zeta'(2,a)$ and the minus sign are essential.

For a function with an endpoint expansion in powers of $x$ and
$\log x$, $\FP\int_0^1$ means the constant coefficient of the cutoff
integral $\int_\varepsilon^1$ as $\varepsilon\downarrow0$, in the
coordinate $x$ itself. We never replace this prescription silently
by a spectral Laurent constant or by a nonlinear cutoff coordinate.

\section{The cubic digamma finite part and the first harmonic coefficient}
\label{bridge:sec}

In this section write $\eta=\eta(1)$ and $Z_2(s,1)=Z_2(s,1;1)$,
with the conventions fixed in Section~\ref{sec:scope}. In particular,
$B_2=(\gamma^2-\zeta(2))/2$ and
\[
 Z_2(1+u,1)=u^{-2}+\gamma u^{-1}+B_2+\eta u+O(u^2).
\]
The inclusive harmonic coefficient is related by
$\eta=-\gamma_H(1,1)-\zeta'(2)$ \cite{Kargin,IncomingOrdered}.

\begin{theorem}[Cubic--harmonic bridge]\label{bridge:main}
In the unit-coordinate Hadamard finite-part convention,
\begin{align}
 Q:=\operatorname{FP}\int_0^1\psi(x)^3\,dx
 &=3\zeta(2)-6\eta-6\gamma\gamma_1,\label{bridge:eta}\\
 &=3\zeta(2)+6\gamma_H(1,1)+6\zeta'(2)-6\gamma\gamma_1.
 \label{bridge:harmonic}
\end{align}
An absolutely convergent arithmetic representation is
\begin{equation}\label{bridge:arithmetic}
 Q=3\zeta(2)+6\gamma_2+
  6\sum_{n=1}^{\infty}
   \frac{\bigl(H_{n-1}-\gamma-\log n\bigr)\log n}{n}.
\end{equation}
Consequently the cubic finite part and the first strict harmonic
coefficient are equivalent modulo ordinary zeta and Stieltjes constants.
This theorem does not assert that the harmonic coefficient is a finite
combination of ordinary zeta constants.
\end{theorem}

\subsection{A global partial-fraction identity}

For $n\ge0$ put
\[
 b_n=H_n-\gamma,\qquad
 a_n=\zeta(2)+H_n^{(2)}-b_n^2,
 \qquad H_n^{(r)}=\sum_{j=1}^n j^{-r},
\]
with $H_0^{(r)}=0$.  Set
\[
 c_0=-\gamma^3+6\gamma\zeta(2)-3\zeta(3).
\]

\begin{lemma}[Normally convergent cubic partial fractions]
\label{bridge:ML}
For $z\notin\{0,-1,-2,\ldots\}$,
\begin{align}
 \psi(z)^3={}&-\frac1{z^3}-\frac{3\gamma}{z^2}
       +\frac{3(\zeta(2)-\gamma^2)}z+c_0\notag\\
 &+\sum_{n=1}^{\infty}\bigg\{
      -\bigl((z+n)^{-3}-n^{-3}\bigr)
      +3b_n\bigl((z+n)^{-2}-n^{-2}\bigr)\notag\\
 &\hspace{42mm}+3a_n\bigl((z+n)^{-1}-n^{-1}\bigr)\bigg\}.
 \label{bridge:MLformula}
\end{align}
The series converges normally on every compact subset avoiding the poles.
\end{lemma}

\begin{proof}
At $z=-n+\varepsilon$, the recurrence of $\psi$ gives
\[
 \psi(-n+\varepsilon)
   =-\varepsilon^{-1}+b_n+
      \bigl(\zeta(2)+H_n^{(2)}\bigr)\varepsilon
      +\bigl(H_n^{(3)}-\zeta(3)\bigr)\varepsilon^2
      +O(\varepsilon^3).
\]
The cubic principal part at $-n$ is therefore
\[
 P_n(z)=-(z+n)^{-3}+3b_n(z+n)^{-2}+3a_n(z+n)^{-1}.
\]
The constant coefficient at zero is $c_0$.
It is important to justify the entire remainder in a Mittag--Leffler
argument, rather than infer it from the principal parts alone.
Let $f(z)=\psi(z)^3-P_0(z)$, which is regular at zero, and integrate
\[
 \frac{z f(w)}{w(w-z)}
\]
over the circles $|w|=N+\tfrac12$, for fixed $z$ and large $N$.
On these circles $f(w)=O(\log^3 N)$ uniformly.  In the right
half-plane this follows from the usual digamma asymptotic \cite[Sections~5.7, 5.11]{DLMFgamma}.  In the
left half-plane use
\[
 \psi(w)=\psi(1-w)-\pi\cot(\pi w).
\]
The cotangent is uniformly bounded: when $|\Im w|\ge1/4$ this follows
from its elementary exponential formula; in the remaining part of the
circle, the real coordinate stays a fixed positive distance from the
integers.  Thus the contour integral is
$O(\log^3 N/N)$ and tends to zero.
Its residues at $z$ and $0$ are $f(z)$ and $-f(0)$.
The residue at $-n$ is $-\bigl(P_n(z)-P_n(0)\bigr)$, either by a
direct calculation or by applying the same two-point Cauchy kernel to
each principal part.  The residue theorem proves
$f(z)-f(0)=\sum_{n\ge1}(P_n(z)-P_n(0))$.
Finally $b_n=O(\log n)$ and $a_n=O(\log^2 n)$, while on compact sets
$(z+n)^{-1}-n^{-1}=O(n^{-2})$.  These estimates prove normal
convergence and justify the asserted identity.
\end{proof}

\subsection{Reduction to one logarithmic harmonic series}

Integrating the normally convergent remainder of
\eqref{bridge:MLformula} term by term yields
\begin{align}
 Q={}&\frac12+3\gamma+c_0
  +\sum_{n\ge1}\bigg\{
    \frac1{n^3}+\frac1{2(n+1)^2}-\frac1{2n^2}
     -\frac{3b_n}{n^2(n+1)}\notag\\
 &\hspace{47mm}
     +3a_n\left(\log\left(1+\frac1n\right)-\frac1n\right)
     \bigg\}.
 \label{bridge:integratedML}
\end{align}
Here the finite parts of $-x^{-3}$ and $-3\gamma x^{-2}$ are
$1/2$ and $3\gamma$, and the unit-coordinate finite part of
$x^{-1}$ is zero.  The elementary identities
\[
 \sum_{n\ge1}\frac{H_n}{n(n+1)}=\zeta(2),\qquad
 \sum_{n\ge1}\frac{H_n}{n^2}=2\zeta(3)
\]
give
\[
 \sum_{n\ge1}\frac{b_n}{n^2(n+1)}
   =2\zeta(3)-\zeta(2)-\gamma\bigl(\zeta(2)-1\bigr).
\]
For completeness, the second Euler sum follows from
\[
 \sum_{n\ge1}\frac{H_n}{n^2}
 =\int_0^1\frac{\log x\log(1-x)}{x(1-x)}\,dx
 =2\int_0^1\frac{\log x\log(1-x)}x\,dx
 =2\zeta(3),
\]
using the harmonic generating function, symmetry, and an absolutely
convergent geometric-series integration.
Consequently \eqref{bridge:integratedML} becomes
\begin{equation}\label{bridge:quadratic-log}
 Q=-\gamma^3+9\gamma\zeta(2)+3\zeta(2)-8\zeta(3)
  +3\sum_{n\ge1}a_n
       \left(\log\left(1+\frac1n\right)-\frac1n\right).
\end{equation}

The crucial cancellation is the exact difference
\[
 a_n-a_{n-1}=-\frac{2(H_{n-1}-\gamma)}n.
\]
Write
$L_N=\sum_{n=1}^N(H_{n-1}-\gamma)\log n/n$ and
$S_N=\sum_{n=1}^N H_n/n^2$.  Finite summation by parts gives
\begin{align*}
 \sum_{n=1}^N a_n\log\left(1+\frac1n\right)
 &=a_N\log(N+1)+2L_N,\\
 \sum_{n=1}^N\frac{a_n}{n}
 &=(\zeta(2)-\gamma^2)H_N+
      \gamma(H_N^2+H_N^{(2)})-\frac{H_N^3}{3}
       +H_NH_N^{(2)}-2S_N+\frac43H_N^{(3)}.
\end{align*}
The second identity follows by summing the finite differences of
$H_n^2$, $H_n^3$, and $H_nH_n^{(2)}$; it has no limiting
interchange hidden in it.  Since
\[
 H_N-\gamma=\log(N+1)+O(N^{-1}),\quad
 H_N^{(2)}=\zeta(2)+O(N^{-1}),\quad S_N\longrightarrow2\zeta(3),
\]
substitution in \eqref{bridge:quadratic-log}, with all polynomial
terms retained before the limit, gives
\begin{equation}\label{bridge:limit}
 Q=3\zeta(2)+6\lim_{N\to\infty}
       \left(L_N-\frac{\log^3 N}{3}\right).
\end{equation}
Now split $H_{n-1}-\gamma=\log n+\psi(n)-\log n$ and use the
standard defining limit
\[
 \gamma_2=\lim_{N\to\infty}
   \left(\sum_{n=1}^N\frac{\log^2 n}{n}-\frac{\log^3 N}{3}\right).
\]
Since $\psi(n)-\log n=O(n^{-1})$, the remaining series is absolutely
convergent.  This proves \eqref{bridge:arithmetic}.

Define the holomorphic Dirichlet deformation
\[
 D(s)=\sum_{n\ge1}\frac{\psi(n)-\log n}{n^s},\qquad\Re s>0.
\]
Normal convergence, including every fixed derivative on compact
subsets, follows from $\psi(n)-\log n=O(n^{-1})$.
In $\Re s>1$ we have the exact identity
\[
 D(s)=Z_2(s,1)-\gamma\zeta(s)+\zeta'(s).
\]
Comparing Laurent coefficients at $s=1$ gives
$D'(1)=\eta+\gamma\gamma_1+\gamma_2$.  Formula
\eqref{bridge:arithmetic} says
$Q=3\zeta(2)+6\gamma_2-6D'(1)$, proving
\eqref{bridge:eta}.  The inclusive-to-strict conversion gives
\eqref{bridge:harmonic} and completes the proof of
Theorem~\ref{bridge:main}.

\subsection{A normalized complex antiderivative}

\begin{corollary}[Cubic digamma primitive]\label{bridge:primitive}
On $\Re z>0$, with the principal logarithm, the function
\begin{align}
 \mathcal P_3(z)={}&\frac1{2z^2}+\frac{3\gamma}z
       +3(\zeta(2)-\gamma^2)\log z+c_0z\notag\\
 &+\sum_{n=1}^{\infty}\bigg\{
      \frac1{2(n+z)^2}-\frac1{2n^2}+\frac{z}{n^3}
      +3b_n\left(\frac1n-\frac1{n+z}-\frac z{n^2}\right)
        \notag\\
 &\hspace{41mm}
      +3a_n\left(\log\left(1+\frac zn\right)-\frac zn\right)
        \bigg\}
 \label{bridge:primitiveformula}
\end{align}
satisfies $\mathcal P_3'(z)=\psi(z)^3$.  Its asymptotic expansion at
zero has constant term zero.  In particular, for every real $b>0$,
\[
 \operatorname{FP}\int_0^b\psi(x)^3\,dx=\mathcal P_3(b),
 \qquad \mathcal P_3(1)=3\zeta(2)-6\eta-6\gamma\gamma_1.
\]
The series and its differentiated series converge normally on compact
subsets of the right half-plane.  Higher derivatives therefore give
identities for derivatives of $\psi^3$ in the same domain.
\end{corollary}

\begin{proof}
Each brace is the primitive, vanishing at $z=0$, of the corresponding
brace in \eqref{bridge:MLformula}.  On a fixed compact set,
the logarithmic difference is $O(n^{-2})$, the quadratic-pole
primitive is $O(n^{-3})$, and the cubic-pole primitive is $O(n^{-4})$.
Multiplication by $a_n=O(\log^2 n)$ or $b_n=O(\log n)$ preserves
normal convergence.  Differentiation recovers the previous lemma.
All series terms vanish at zero, so the displayed singular and logarithmic
terms leave no constant term.  This proves the stated normalization.
\end{proof}

\begin{corollary}[A polygamma--harmonic hierarchy]\label{bridge:polygamma}
For every integer $r\ge1$ and $z\notin\{0,-1,-2,\ldots\}$,
\begin{align}
 &\sum_{r_1+r_2+r_3=r}\frac{r!}{r_1!r_2!r_3!}
       \psi^{(r_1)}(z)\psi^{(r_2)}(z)\psi^{(r_3)}(z)\notag\\
 &\quad=(-1)^r r!\sum_{n=0}^{\infty}
   \left\{-\frac{(r+1)(r+2)}{2(n+z)^{r+3}}
       +\frac{3(r+1)b_n}{(n+z)^{r+2}}
       +\frac{3a_n}{(n+z)^{r+1}}\right\}.
 \label{bridge:allpolygamma}
\end{align}
Here $\psi^{(0)}=\psi$. The series converges normally away from
the poles. At $z=1$ the left side is a finite polynomial in
$\gamma,\zeta(2),\ldots,\zeta(r+1)$, using
$\psi^{(k)}(1)=(-1)^{k+1}k!\zeta(k+1)$ for $k\ge1$.
\end{corollary}
\begin{proof}
Differentiate \eqref{bridge:MLformula} $r$ times. The constant
subtractions disappear and the derivative series is bounded on each
compact set by a constant times
$\sum n^{-r-1}(1+\log^2n)$. The product rule gives the left
side. The $n=0$ summand includes the principal part at zero since
$b_0=-\gamma$ and $a_0=\zeta(2)-\gamma^2$. The special values
follow from the classical Hurwitz--polygamma identity
\cite[Section~25.11]{DLMFhurwitz}.
\end{proof}

\subsection{A short certified evaluation route}

For integers $N\ge2$ and $K\ge0$ define
\begin{align*}
 S_{N,K}={}&\sum_{n=1}^{N-1}
  \frac{(H_{n-1}-\gamma-\log n)\log n}{n}
   +\frac12\zeta'(2,N)
   +\sum_{k=1}^K\frac{B_{2k}}{2k}\zeta'(2k+1,N).
\end{align*}
The Binet formula for $\psi$ \cite[Section~5.9]{DLMFgamma} and the finite geometric expansion of
$(n^2+t^2)^{-1}$ give the rigorous truncation estimate
\begin{equation}\label{bridge:tail-bound}
 \left|S_{N,K}-\sum_{n\ge1}
  \frac{(\psi(n)-\log n)\log n}{n}\right|
 \le \frac{|B_{2K+2}|}{2K+2}\bigl(-\zeta'(2K+3,N)\bigr).
\end{equation}
Indeed the digamma asymptotic remainder is bounded by the magnitude of
its first omitted Bernoulli term for positive real arguments, and the
logarithmic tail is exactly a Hurwitz-zeta derivative.
The supplied Arb script evaluates every displayed quantity with ball
arithmetic and adds this analytic tail bound.  This certifies the
arithmetic identity numerically without replacing its proof.

\paragraph{Status.}
The harmonic zeta function, its Stieltjes coefficients, and the
Mittag--Leffler method are classical.  The cubic--harmonic bridge and
its coordinate consequence were derived here and were not found in the
inspected repository reports.  No worldwide priority claim is made.

\section{The diagonal Tornheim derivative in harmonic coordinates}
\label{sec:diagonal}

The incoming coincident report relates the digamma cube to a third
diagonal Tornheim derivative. We give the relevant derivation, with
its finite normalization visible, and then insert the new bridge.
This separates the earlier spectral identity from the additional
harmonic reduction proved here.

Write
\[
 \omega(t)=\T(t,t,t),\qquad \Omega_3=\omega'''(0),\qquad
 E(t)=\Gamma(1+t)e^{-Lt},\qquad
 h(t)=\zeta(1-t)+\frac1t.
\]
Here $\T$ is the continuation of
$\sum_{m,n\ge1}m^{-A}n^{-B}(m+n)^{-C}$, and the diagonal is
restricted before taking the limit at zero. The third derivative is
one-variable data; joint holomorphy at the origin is not assumed.

\subsection{The completed cubic generating function}

Set
\[
 g(t,x)=\zeta(1-t,x)+\frac1t
       =\sum_{m\ge0}\gamma_m(x)\frac{t^m}{m!}.
\]
Near $x=0$ its exact Hurwitz splitting yields
\begin{equation}\label{diag:local}
 g(t,x)=x^{t-1}+h(t)-(1-t)\zeta(2-t)x+O(x^2).
\end{equation}
After cubing, the only local monomials whose Mellin denominators
vanish at $t=0$ are
\[
 3h(t)^2x^{t-1},\qquad
 -3(1-t)\zeta(2-t)x^{2t-1}.
\]
Their complete spectral integrals are $3h(t)^2/t$ and
$-3(1-t)\zeta(2-t)/(2t)$. The complete numerators, rather than
just their values at zero, have to be removed to obtain the
unit-coordinate finite parts.

For $\Re t>1$, the absolutely convergent Hurwitz Fourier expansion
has coefficients
\[
 \widehat{\zeta(1-t,\cdot)}(n)
 =\Gamma(t)(2\pi|n|)^{-t}
             e^{-\pi i t\operatorname{sgn}(n)/2},\qquad n\ne0.
\]
The constant coefficient vanishes. In the product of three series,
the condition $n_1+n_2+n_3=0$ has six sign cones; on each, the
largest magnitude is the sum of the other two. Consequently
\begin{align}
 \int_0^1\zeta(1-t,x)^3\dd x
 &=6\Gamma(t)^3(2\pi)^{-3t}\cos(\pi t/2)\omega(t),\notag\\
 \int_0^1\zeta(1-t,x)^2\dd x
 &=2\Gamma(t)^2(2\pi)^{-2t}\zeta(2t).
 \label{diag:Fourier}
\end{align}
The classical Hurwitz--Tornheim connection is recorded, with its
literature attribution, in the coincident report
\cite{IncomingCoincident}. Finite endpoint Taylor subtraction
continues both sides meromorphically.

\begin{lemma}[Completed diagonal kernel]\label{lem:diag-completed}
The expression
\begin{align}
 F(t)={}&\frac{6E(t)^3\cos(\pi t/2)\omega(t)
                         +6E(t)^2\zeta(2t)+1}{t^3}\notag\\
 &-\frac{3h(t)^2}{t}
           +\frac{3(1-t)\zeta(2-t)}{2t}
 \label{diag:completed}
\end{align}
is holomorphic near zero and satisfies $F(0)=-Q$.
\end{lemma}

\begin{proof}
Expand $g^3$ and use \eqref{diag:Fourier} and the vanishing
one-factor integral. This gives the fraction on the first line as
the meromorphic integral of $g^3$. To justify its completion,
split the $x$ integral at any $b\in(0,1)$. Subtract finitely many
Taylor monomials of the smooth factor in \eqref{diag:local} on
$(0,b)$. The remainder and all fixed parameter derivatives are
integrable uniformly in a small $t$ disk. A local monomial with
exponent $kt-1$ has spectral integral $b^{kt}/(kt)$, while the
generating function of its coordinate finite parts is
$(b^{kt}-1)/(kt)$. Their difference is exactly $1/(kt)$.
All nonresonant denominators remain nonzero at zero. The two
resonant terms displayed above therefore give precisely the second
line of \eqref{diag:completed}. The resulting Taylor coefficients
are the finite parts of those of $g^3$. At zero,
$g(0,x)=\gamma_0(x)=-\psi(x)$, hence $F(0)=-Q$.
\end{proof}

Solving \eqref{diag:completed} for $\omega$ also proves its
holomorphy at zero, since the coefficient $E(t)^3\cos(\pi t/2)$
is nonzero there. This argument does not infer a joint regularity
statement from a ray restriction.

\subsection{An explicit reduction of the third diagonal derivative}

Use the classical values $\zeta(0)=-1/2$, $\zeta'(0)=-L/2$ and
\[
 E(t)=\exp\left[-(\gamma+L)t+\frac{\zeta(2)}2t^2
                              -\frac{\zeta(3)}3t^3+O(t^4)\right].
\]
The coefficients of $t^{-3},t^{-2},t^{-1}$ in
\eqref{diag:completed} give
\begin{equation}\label{diag:lower}
 \omega(0)=\frac13,\qquad \omega'(0)=L,
 \qquad \omega''(0)=L^2-4\zeta''(0).
\end{equation}
For clarity, after the first two substitutions the coefficient of
$t^2$ in the numerator is
\[
 3\omega''(0)+3\gamma^2-3L^2-\tfrac32\zeta(2)+12\zeta''(0),
\]
and the remaining two terms contribute
$-3\gamma^2+3\zeta(2)/2$ to the residue.
The numerator coefficient at $t^3$, after \eqref{diag:lower}, is
\[
 \Omega_3+8\zeta'''(0)+12(\gamma+L)\zeta''(0)
 +L^3+6\gamma L^2-5\gamma^3+\tfrac32\gamma\zeta(2).
\]
Since $h(t)=\gamma+\gamma_1t+O(t^2)$, the other constant
terms are $-6\gamma\gamma_1-3(\zeta(2)+\zeta'(2))/2$.
Thus the earlier cubic identity is recovered:
\begin{align}
 Q={}&-\Omega_3-8\zeta'''(0)-12(\gamma+L)\zeta''(0)
            -L^3-6\gamma L^2+5\gamma^3\notag\\
 &+6\gamma\gamma_1-\frac32\gamma\zeta(2)
                       +\frac32\bigl(\zeta(2)+\zeta'(2)\bigr).
 \label{diag:old-cube}
\end{align}

\begin{theorem}[Harmonic coordinate for the cubic diagonal]
\label{thm:diag-harmonic}
The third derivative of the specified diagonal continuation is
\begin{align}
 \Omega_3={}&6\eta(1)-8\zeta'''(0)
 -12(\gamma+L)\zeta''(0)-L^3-6\gamma L^2+5\gamma^3\notag\\
 &+12\gamma\gamma_1-\frac32\gamma\zeta(2)
                     -\frac32\zeta(2)+\frac32\zeta'(2).
 \label{bridge:tornheim}
\end{align}
In particular, \eqref{bridge:arithmetic} gives an absolutely
convergent generalized-harmonic evaluation of this derivative.
\end{theorem}
\begin{proof}
Substitute \eqref{bridge:eta} into \eqref{diag:old-cube} and solve
for $\Omega_3$.
\end{proof}

The independently certified arithmetic evaluation in the package gives
\begin{align*}
 Q&=1.96626273439153434066437533841163142382586897649\ldots,\\
 \eta(1)&=0.53678702435841856750360788434881167385105703307\ldots,\\
 \Omega_3&=86.6578666011000736529032768587621836100797302472\ldots.
\end{align*}
The first harmonic coefficient is therefore sufficient for the
diagonal cubic problem. A finite expression for that coefficient in
a specified algebra of ordinary zeta, Gamma and Stieltjes constants
remains an arithmetic research question.

\section{The complete cubic directional layer of Tornheim zeta}
\label{sec:rays}

Define the uncolored Mordell--Tornheim function in its absolute
convergence region by
\begin{equation}\label{ray:T}
 \T(A,B,C)=\sum_{m,n\ge1}\frac1{m^A n^B(m+n)^C}.
\end{equation}
For instance, sufficiently large positive real parts suffice.
Its meromorphic continuation is symmetric in $A,B$, but is not
jointly holomorphic at the origin. We form a restriction away from
its polar hyperplanes before taking a one-variable limit.

Let
\[
 W_{a,b,c}(t)=\T(at,bt,ct),\qquad
 \Omega_3=\left.\frac{\mathrm d^3}{\mathrm dt^3}\T(t,t,t)\right|_{t=0}.
\]
The two admissibility conditions are
\begin{equation}\label{ray:admissible}
 (a+c)(b+c)\ne0.
\end{equation}
They allow complex and unequal slopes, including negative slopes.

\subsection{Statement of the formula}

Put
\begin{align}
 \sigma&=a+b+c,\notag\\
 \mathcal B&=\frac{c(c^2+ab-a^2-b^2)}{(a+c)(b+c)},\label{ray:B}\\
 \Delta&=a^2+b^2-c(a+b),\notag\\
 J_3&=W_{1,0,2}'''(0),\notag\\
 \kappa&=J_3+\frac{35}{2}\zeta'''(0)
                 +6L\zeta''(0)-9\zeta'''(-1).\label{ray:kappa}
\end{align}
The constant $J_3$ is a fixed Euler double-zeta directional derivative:
$\T(t,0,2t)=Z_2(2t,t;1)$. An independent convergent integral for
$\kappa$ is given below.

\begin{theorem}[All admissible cubic directions]\label{thm:cubic-rays}
Under \eqref{ray:admissible}, $W_{a,b,c}$ is holomorphic at zero and
\begin{align}
 W_{a,b,c}'''(0)={}&abc\,\Omega_3-\frac{c\Delta}{2}\kappa\notag\\
 &+\left(8abc-\frac{(a+c)^3+(b+c)^3}{2}\right)\zeta'''(0)\notag\\
 &-\frac32\bigl((a+b)(ab+c^2)-4abc\bigr)L\zeta''(0)
       +\sigma^2\mathcal B\,\zeta'''(-1).
 \label{ray:main}
\end{align}
Thus every third directional derivative belongs to the span of
$\Omega_3,\kappa,\zeta'''(0),L\zeta''(0),\zeta'''(-1)$ over
the field of rational functions of the slopes. By the cubic bridge,
$\Omega_3$ can be replaced by $\eta(1)$ and explicit ordinary zeta
and Stieltjes data.
\end{theorem}

Every coefficient is homogeneous of degree three. This agrees with
$W_{ra,rb,rc}(t)=W_{a,b,c}(rt)$. The formula contains no surviving
Euler-constant or Gamma-expansion coefficients in addition to those
already packaged into $\Omega_3$ and $\kappa$.

\subsection{A holomorphic remainder with exactly two visible polar divisors}

For $x>0$ set $F(A,B;x)=\Li_A(e^{-x})\Li_B(e^{-x})$.
Jonqui\`ere's formula gives, locally uniformly for $A$ near zero,
\begin{equation}\label{ray:Jonquiere}
 \Li_A(e^{-x})=\Gamma(1-A)x^{A-1}
       +\sum_{j\ge0}\zeta(A-j)\frac{(-x)^j}{j!},\qquad 0<x<2\pi.
\end{equation}
Subtract the six local terms
\begin{align}
 S(A,B;x)={}&\Gamma(1-A)\Gamma(1-B)x^{A+B-2}
 +\Gamma(1-A)\zeta(B)x^{A-1}\notag\\
 &+\Gamma(1-B)\zeta(A)x^{B-1}
 -\Gamma(1-A)\zeta(B-1)x^A\notag\\
 &-\Gamma(1-B)\zeta(A-1)x^B+\zeta(A)\zeta(B).
 \label{ray:S}
\end{align}
Define
\begin{align}
 H(A,B,C)={}&\int_0^1x^{C-1}\bigl(F(A,B;x)-S(A,B;x)\bigr)\dd x
       +\int_1^\infty x^{C-1}F(A,B;x)\dd x\notag\\
 &+\frac{\Gamma(1-A)\Gamma(1-B)}{A+B+C-2}
 +\frac{\Gamma(1-A)\zeta(B)}{A+C-1}
 +\frac{\Gamma(1-B)\zeta(A)}{B+C-1}.
 \label{ray:H}
\end{align}
All denominators on the last line are nonzero near the origin.
On a sufficiently small closed parameter polydisc, the first integrand
and every fixed parameter derivative are bounded at zero by
$x^{-\delta}(1+|\log x|^N)$ for a $\delta<1$; the second integral
has exponential decay. Thus $H$ is jointly holomorphic and symmetric
in $A,B$. Termwise integration and meromorphic continuation yield
\begin{equation}\label{ray:decomposition}
 \begin{split}
 \T(A,B,C)=\frac1{\Gamma(1+C)}\biggl\{&\zeta(A)\zeta(B)
 -\frac{C\Gamma(1-A)\zeta(B-1)}{A+C}\\
 &-\frac{C\Gamma(1-B)\zeta(A-1)}{B+C}
 +CH(A,B,C)\biggr\}.
 \end{split}
\end{equation}
This also proves removability on every admissible ray. The six-term
subtraction and the decomposition are retained from the incoming
coincident report \cite{IncomingCoincident}; the cubic determination
below is the additional step.

\subsection{Exact slices determine the constant and linear remainder}

Two elementary restrictions provide all the lower coefficients:
\begin{align}
 \T(0,0,t)&=\zeta(t-1)-\zeta(t),\label{ray:axis}\\
 \T(t,0,t)&=\frac{\zeta(t)^2-\zeta(2t)}2.\label{ray:symmetric-slice}
\end{align}
For the first, count positive pairs with $m+n=k$. For the second,
use $\T(A,0,C)=Z_2(C,A;1)$ and the exact stuffle product
\begin{equation}\label{ray:stuffle-slice}
 \T(A,0,C)+\T(C,0,A)=\zeta(A)\zeta(C)-\zeta(A+C).
\end{equation}
Both identities are first proved in domains of absolute convergence
and then extended meromorphically. A fixed $C=0$ substitution alone
would not determine the normal derivative of $CH$; this is why
\eqref{ray:symmetric-slice} supplies real additional information.

Write
\begin{align}
 H(A,B,C)={}&h_0+h_A(A+B)+h_C\,C\notag\\
 &+h_{20}(A^2+B^2)+h_{11}AB+h_{10c}(A+B)C+h_{02}C^2
 +O\bigl((|A|+|B|+|C|)^3\bigr).
 \label{ray:H-Taylor}
\end{align}
Here $h_C$ denotes the coefficient of the linear term in $C$;
$h_{02}$ is its quadratic coefficient. Direct comparison in
\eqref{ray:axis}--\eqref{ray:symmetric-slice} gives
\begin{align}
 h_0={}&\zeta'(-1)+\frac L2-\frac{5\gamma}{12},\label{ray:h0}\\
 h_C={}&\frac{\zeta''(-1)-\zeta''(0)}2
       -\gamma\left(\zeta'(-1)+\frac L2\right)
       +\frac5{24}(\gamma^2+\zeta(2)),\label{ray:hC}\\
 h_A={}&\frac{L^2}{8}-\frac{\zeta''(0)}2+
       \gamma\zeta'(-1)-\frac{\gamma L}{4}
       -\frac{\gamma^2+\zeta(2)}{24}.
 \label{ray:hA}
\end{align}
For example, the axis gives
$\Gamma(1+t)(\zeta(t-1)-\zeta(t))=5/12+tH(0,0,t)$,
which proves the first two lines. On the symmetric slice the
coefficient of $t^2$ then gives the third line.

Substitution now proves the complete lower layer
\begin{align}
 W(0)&=\frac14+\frac{c}{12(a+c)}+\frac{c}{12(b+c)},\notag\\
 W'(0)&=\frac{a+b+2c}{4}L+\mathcal B\zeta'(-1),\notag\\
 W''(0)&=\frac{2ab+c(a+b)}4L^2
 -\frac{a^2+b^2+2c(a+b)+2c^2}{2}\zeta''(0)
 +\sigma\mathcal B\zeta''(-1).
 \label{ray:lower}
\end{align}
These agree with the earlier report. In particular,
 $\omega(0)=1/3$, $\omega'(0)=L$, and
$\omega''(0)=L^2-4\zeta''(0)$ follow here without using the
unknown cubic moment.

\subsection{Proof of the cubic formula}

At quadratic degree, the symmetric slice fixes
$h_{20}+h_{10c}$, not the two terms separately. This distinction is
the main structural point. The axis fixes $h_{02}$; the diagonal
then fixes $h_{11}$, leaving exactly one more coefficient, determined
by $J_3$.

Here are explicit finite coefficient formulas which also make the
algebra completely reproducible. Set $V=h_{20}+h_{10c}$ and define
the cubic Taylor polynomials
\begin{align*}
 O(t)&=\frac13+Lt+\frac{L^2-4\zeta''(0)}2t^2
                         +\frac{\Omega_3}{6}t^3,\\
 D(t)&=\frac7{18}+\left(\frac54L+\zeta'(-1)\right)t
       +\left(\frac{L^2}{4}-\frac{13}{4}\zeta''(0)
                               +\frac32\zeta''(-1)\right)t^2
       +\frac{J_3}{6}t^3.
\end{align*}
Comparison on the four slices gives
\begin{align}
 h_{02}&=[t^3]\Gamma(1+t)(\zeta(t-1)-\zeta(t)),\notag\\
 V+h_{02}&=[t^3]\left\{
 \frac{\Gamma(1+t)}2(\zeta(t)^2-\zeta(2t))
 +\frac{\zeta(t)}2-\frac{\Gamma(1-t)}{24}+\zeta(t-1)\right\},\notag\\
 h_{11}+2V+h_{02}&=[t^3]\left\{
 \Gamma(1+t)O(t)-\zeta(t)^2+\Gamma(1-t)\zeta(t-1)\right\},\notag\\
 -2h_{20}+4V+8h_{02}&=[t^3]\left\{
 \Gamma(1+2t)D(t)+\frac{\zeta(t)}2
                      -\frac{\Gamma(1-t)}{18}+\zeta(t-1)\right\}.
 \label{ray:finite-system}
\end{align}
The coefficient system is triangular with nonzero diagonal. Thus
it determines all four quadratic coefficients for any assigned
$\Omega_3,J_3$.

Use
\[
 \Gamma(1\pm t)=1\mp\gamma t+
 \frac{\gamma^2+\zeta(2)}2t^2
 \mp\frac{\gamma^3+3\gamma\zeta(2)+2\zeta(3)}6t^3+O(t^4)
\]
and the ordinary Taylor expansions of $\zeta$ at $0$ and $-1$.
Substitute \eqref{ray:finite-system} into
\eqref{ray:decomposition}. The coefficients of $\Omega_3$ and
$J_3$ are respectively $abc$ and $-c\Delta/2$.
The remaining nonzero coefficients are
\begin{align*}
 [\zeta'''(0)]\colon\;&8abc-\frac{(a+c)^3+(b+c)^3}{2}
                                      -\frac{35}{4}c\Delta,\\
 [L\zeta''(0)]\colon\;&-\frac32\bigl((a+b)(ab+c^2)-4abc\bigr)
                                      -3c\Delta,\\
 [\zeta'''(-1)]\colon\;&\sigma^2\mathcal B+\frac92c\Delta.
\end{align*}
All coefficients involving $\gamma,\zeta(2),\zeta(3)$ or
$\zeta'(-1),\zeta''(-1)$ cancel. Absorbing the three multiples
of $c\Delta$ into \eqref{ray:kappa} proves
Theorem~\ref{thm:cubic-rays}. This is a finite algebraic proof;
the delivered symbolic certificate independently verifies the exact
rational identity with all constants treated as formal symbols.

\subsection{An independent convergent integral for the second coordinate}

Define
\begin{align}
 Q_2(x)&=\left.\partial_s^2\Li_s(e^{-x})\right|_{s=0}
               =\sum_{n\ge1}(\log n)^2e^{-nx},\notag\\
 P(x)&=(\log x+\gamma)^2+\zeta(2),\notag\\
 S_2(x)&=\left(x^{-2}-\frac1{2x}+\frac1{12}\right)P(x)
             +\zeta''(0)\left(x^{-1}-\frac12\right)-\zeta''(-1),\notag\\
 I_2&=\int_0^1\left(\frac{Q_2(x)}{e^x-1}-S_2(x)\right)\frac{\dd x}{x}
       +\int_1^\infty\frac{Q_2(x)}{x(e^x-1)}\dd x.
 \label{ray:I2}
\end{align}
Both integrals converge ordinarily: the first integrand is
$O(1+\log^2x)$; the second is $O(e^{-3x}/x)$.

\begin{proposition}[Mellin coordinate]\label{prop:kappa-integral}
The quadratic coefficient $h_{20}$ satisfies
\begin{equation}\label{ray:h20-integral}
 h_{20}=\frac\gamma4+\frac38-\frac{\zeta''(0)}2+\frac{I_2}{2},
\end{equation}
and
\begin{align}
 \kappa={}&-12h_{20}-\frac{\gamma^3}{6}
       -\frac{\gamma\zeta(2)}2+3\gamma\zeta''(0)
       +6\gamma\zeta''(-1)-\frac{\zeta(3)}3
       -3\zeta'''(0)-2\zeta'''(-1).
 \label{ray:kappa-h20}
\end{align}
\end{proposition}
\begin{proof}
Differentiate \eqref{ray:H} twice in $A$ at $(0,0,0)$.
The product derivative is $Q_2(x)/(e^x-1)$, and the derivative of
the six subtractions is $S_2(x)$. Differentiation under the subtracted
integrals follows from the local bounds used above. The rational
last line of \eqref{ray:H}, on $B=C=0$, is
\[
 \Gamma(1-A)\left(\frac1{A-2}-\frac1{2(A-1)}\right)-\zeta(A).
\]
Its coefficient of $A^2$ is
$\gamma/4+3/8-\zeta''(0)/2$, proving
\eqref{ray:h20-integral}. Expanding the last line of
\eqref{ray:finite-system} and using \eqref{ray:kappa} gives
\eqref{ray:kappa-h20}.
\end{proof}

The independent integral evaluation gives
\begin{align*}
 I_2&=-0.2860481161626797704387366252971602\ldots,\\
 \kappa&=-3.0226121641572464811701683702659737\ldots.
\end{align*}
These displayed digits are numerical diagnostics, not an arithmetic
characterization. The definition \eqref{ray:I2} supplies such a
characterization without referring back to an unknown generic ray.

\subsection{Coordinate cancellation and exact examples}

In the formal coordinate system of Theorem~\ref{thm:cubic-rays},
the $\kappa$ term disappears exactly when
\begin{equation}\label{ray:conic}
 c\bigl(a^2+b^2-c(a+b)\bigr)=0.
\end{equation}
For positive $a,b$, the choice
$c=(a^2+b^2)/(a+b)$ gives a continuum of asymmetric directions with
only the harmonic coordinate $\eta(1)$ remaining. For fixed real
$c$, the nontrivial locus can equivalently be written
$(a-c/2)^2+(b-c/2)^2=c^2/2$.

The $\Omega_3$ term disappears when $abc=0$. With $c\ne0$, both
displayed nonclassical coefficients vanish only on the axis
$(0,0,c)$ and the rays $(c,0,c)$ or $(0,c,c)$. With $c=0$,
admissibility requires $ab\ne0$ and the restriction equals
$\zeta(at)\zeta(bt)$. These are coefficient-elimination statements,
not assertions that other numerical cancellations are impossible.

Some explicit specializations are
\begin{align}
 W_{1,2,3}'''(0)={}&6\Omega_3+6\kappa
       -\frac{93}{2}\zeta'''(0)-\frac{27}{2}L\zeta''(0)
       +\frac{162}{5}\zeta'''(-1),\label{ray:123}\\
 W_{2,3,1}'''(0)={}&6\Omega_3-4\kappa
       +\frac52\zeta'''(0)-\frac{33}{2}L\zeta''(0)
       -18\zeta'''(-1),\label{ray:231}\\
 W_{1,1,2}'''(0)={}&2\Omega_3+2\kappa
       -11\zeta'''(0)-3L\zeta''(0)+\frac{32}{3}\zeta'''(-1),\notag\\
 W_{0,0,1}'''(0)={}&\zeta'''(-1)-\zeta'''(0),\notag\\
 W_{1,0,1}'''(0)={}&-\frac92\zeta'''(0)-\frac32L\zeta''(0).
\end{align}
The last two also follow directly from the exact slices.

\begin{remark}[Why two coordinates are sufficient, and what this does not prove]
Modulo the already determined constant and linear coefficients of
$H$, the axis and symmetric slice leave two formal quadratic degrees
of freedom. The diagonal and $(1,0,2)$ slices fix them separately.
Their contributions to generic cubic rays are independent polynomials
$abc$ and $c\Delta$. This proves uniqueness in the specified formal
reconstruction problem. It does not prove that $\eta(1)$ and $\kappa$
are independent over any arithmetic constant field.
\end{remark}

\section{Transverse curved directions and a sharp fourth-jet correction}
\label{sec:curves}

The cubic coordinate formula survives a nonlinear change of spectral
path. The path contributes an explicitly classical correction, while
the coefficients of $\Omega_3$ and $\kappa$ depend only on its tangent.

Using \eqref{ray:h0}--\eqref{ray:hA}, define
\begin{align}
 K(A,B,C)=\frac1{\Gamma(1+C)}\biggl\{
 &\zeta(A)\zeta(B)
 -\frac{C\Gamma(1-A)\zeta(B-1)}{A+C}
 -\frac{C\Gamma(1-B)\zeta(A-1)}{B+C}\notag\\
 &+C\bigl(h_0+h_A(A+B)+h_C\,C\bigr)\biggr\}.
 \label{curve:K}
\end{align}
Although $K$ is meromorphic in three variables, its restriction to
each of the paths in the following theorem is holomorphic.

\begin{theorem}[Cubic transport along a spectral curve]\label{thm:curve}
Let $A(t),B(t),C(t)$ be analytic near zero, vanish at zero, and have
first derivatives $a,b,c$ satisfying \eqref{ray:admissible}.
Write $\mathscr W(t)=\T(A(t),B(t),C(t))$. Then
\begin{equation}\label{curve:main}
 \mathscr W'''(0)=W_{a,b,c}'''(0)+
 6[t^3]\bigl(K(A(t),B(t),C(t))-K(at,bt,ct)\bigr).
\end{equation}
The correction uses only ordinary Gamma and zeta derivatives and
the Taylor coefficients of $A,B,C$ through degree four. In particular,
the coefficients of $\Omega_3$ and $\kappa$ remain $abc$ and
$-c\Delta/2$.
\end{theorem}

\begin{proof}
Write the quadratic homogeneous part of \eqref{ray:H-Taylor} as
$H_2$. Formula \eqref{ray:decomposition} gives
\[
 \T-K=\frac{C}{\Gamma(1+C)}H_2(A,B,C)
            +O\bigl((|A|+|B|+|C|)^4\bigr).
\]
After substitution of the curve, its coefficient of $t^3$ is
$cH_2(a,b,c)$, exactly the coefficient on the tangent ray. Subtracting
the two restrictions proves \eqref{curve:main}.

The only possible need for one additional path coefficient comes
from $C/(A+C)$ and $C/(B+C)$. Remove their common factor $t$.
The new denominators have nonzero constant terms, so the coefficient
of $t^3$ in each ratio uses the original curves only through degree
four. All the remaining compositions use at most degree three.
This proves the finite-jet assertion and holomorphy of the restrictions.
\end{proof}

\subsection{A correction involving only the lower directional layer}

The correction can be made smaller than the Gamma expression
\eqref{curve:K} suggests. Put $v=(a,b,c)$ and write the path as
\[
 (A(t),B(t),C(t))=tv+t^2\alpha+t^3\beta+t^4\delta+O(t^5).
\]
Let $f_j(v)=W_v^{(j)}(0)$ for $j=0,1,2$, explicitly given by
\eqref{ray:lower}. Denote ordinary differentials in the slope
variables by $D$. Then
\begin{align}
 \mathscr W'''(0)-W_v'''(0)={}&
 6Df_0(v)[\delta]+6D^2f_0(v)[\alpha,\beta]
                  +D^3f_0(v)[\alpha,\alpha,\alpha]\notag\\
 &+6Df_1(v)[\beta]+3D^2f_1(v)[\alpha,\alpha]
                  +3Df_2(v)[\alpha].
 \label{curve:differentials}
\end{align}
In particular the correction belongs to the span of
$1,L,L^2,\zeta'(-1),\zeta''(0),\zeta''(-1)$ with rational
functions of the path coefficients. No third zeta derivative or
additional Gamma constant is needed.

To prove \eqref{curve:differentials}, cancel the common factor $t$
in the two polar ratios in \eqref{ray:decomposition}. The function
$(t,v)\mapsto\T(ta,tb,tc)$ is then jointly holomorphic near
$(0,v)$ for admissible $v$. Its Taylor series is
$\sum_{j\ge0}f_j(v)t^j/j!$. Substitute
$v+\alpha t+\beta t^2+\delta t^3+O(t^4)$ for $v$ and take
the coefficient of $t^3$. The ordinary multivariate Taylor formula
gives exactly \eqref{curve:differentials}.

\begin{proposition}[The fourth jet is genuinely needed]\label{prop:sharpcurve}
Fix admissible $a,b,c$ with $c\ne0$. Replacing the straight path
$(at,bt,ct)$ by $(at+\varepsilon t^4,bt,ct)$ changes its cubic
derivative by exactly
\begin{equation}\label{curve:sharp}
 -\frac{c\varepsilon}{2(a+c)^2}.
\end{equation}
Here $\varepsilon$ is the coefficient of $t^4$, not the fourth derivative.
\end{proposition}
\begin{proof}
Only the ratio $C/(A+C)$ in \eqref{curve:K} changes at degree three.
Its multiplier at the origin is $-\zeta(-1)=1/12$. The ratio changes by
$-c\varepsilon t^3/(a+c)^2+O(t^4)$. Multiplication by $3!=6$
gives \eqref{curve:sharp}.
\end{proof}

For example, the curve $(t+\varepsilon t^4,t,t)$ has cubic
derivative $\Omega_3-\varepsilon/8$, even though its first three
path derivatives agree with those of the diagonal.

This theorem concerns transverse curves. If a path lies in $A+C=0$
or $B+C=0$, or is tangent to one of those divisors, a separate
restriction or normal finite-part prescription is needed. The theorem
does not assign a value by substituting a zero denominator into
\eqref{ray:main}.

% Integration requirements: amsmath, amssymb, amsthm; theorem, lemma,
% proposition, corollary, remark, and proof environments.  No unprefixed
% macros are introduced.  Bibliography entries are in mixed_references.tex.
\providecommand{\mxLi}{\operatorname{Li}}

\section{Arbitrary depth with mixed integer inner orders}
\label{mx:section}

Question~11 of the preceding directional report asks for a finite
nested-polylogarithm reduction when the inner integer orders have both
signs.  Theorem~\ref{mx:main} supplies such a reduction while keeping the
outer order complex.  If exactly $q$ inner orders are positive, its
output has depth at most $q+1$, regardless of how many nonpositive
orders occur.  Its coefficients are finite rational operations,
root-of-unity filters, and generalized-harmonic polynomials.  Repeated
poles are included directly, and coincident colors are evaluated by an
explicit polynomial summation rule.

The shuffle algebra of iterated integrals and the reduction of rational
functions times hyperlogarithms are established tools; see
Panzer~\cite[Sections~2.1--2.3]{mx:Panzer}.  Colored depth-two Tornheim
reductions are also classical~\cite{mx:Zhao}.  The contribution here is
the explicit convolution formula with unrestricted outer order, its
depth bound, its coincident-color reduction rule, and the accompanying
meromorphic and normal-derivative prescription.  No claim of literature
priority for the underlying shuffle or rational-function algorithms is
made.

\subsection{Definitions and finite input data}

Let $s_j\in\mathbb Z$, $p_j\in\mathbb Z_{>0}$, and
$0<|X_j|\leq1$.  Define
\begin{equation}\label{mx:Z}
 Z_{\mathbf s,\mathbf p}(C;\mathbf X)
 =\sum_{m_1,\ldots,m_d\geq1}
  \frac{\prod_{j=1}^dX_j^{m_j}m_j^{-s_j}}
       {(p_1m_1+\cdots+p_dm_d)^C}.
\end{equation}
The half-plane
$\Re C>d+\sum_j\max(0,-s_j)$ is sufficient for absolute
convergence.  Every power of a positive integer uses its real
logarithm.  Put
\[
 I_+=\{j:s_j>0\},\qquad I_-=\{j:s_j\leq0\},\qquad
 q=|I_+|,
\]
and assume for the main theorem that both sets are nonempty.  Factor
the ordinary coefficient generating function as
\begin{equation}\label{mx:factor}
 \sum_{n\geq1}r_nt^n=R(t)P(t),\qquad
 R(t)=\prod_{j\in I_-}\mxLi_{s_j}(X_jt^{p_j}),\quad
 P(t)=\prod_{j\in I_+}\mxLi_{s_j}(X_jt^{p_j}).
\end{equation}
Then $Z(C)=\sum_{n\geq1}r_nn^{-C}$ in the stated half-plane.

For each $j\in I_-$ write $a_j=-s_j$.  The identity
\[
 \mxLi_{-a}(z)=\left(z\frac{d}{dz}\right)^a\frac z{1-z}
\]
shows that $R$ is rational and bounded at infinity.  Its finite partial
fractions have the form
\begin{equation}\label{mx:partial}
 R(t)=c_\infty+
       \sum_\alpha\sum_{h=1}^{M_\alpha}
          \frac{c_{\alpha,h}}{(1-\alpha t)^h},\qquad
 M_\alpha=\sum_{\substack{j\in I_-\\\alpha^{p_j}=X_j}}(1-s_j).
\end{equation}
The index $\alpha$ runs over the distinct roots of the indicated
equations.  Pole cancellation may make some coefficients zero.  Even
at repeated roots the coefficient formula is simply
\begin{equation}\label{mx:polecoeff}
 c_{\alpha,h}
 =[u^{M_\alpha-h}]\,u^{M_\alpha}R((1-u)/\alpha).
\end{equation}
The function inside the coefficient extraction is holomorphic at zero.
Because $I_-$ is nonempty, $R(0)=0$, so
\begin{equation}\label{mx:constantcancel}
 c_\infty+\sum_{\alpha,h}c_{\alpha,h}=0.
\end{equation}
Finally set
\begin{equation}\label{mx:e}
 e_{h,r}=[v^r]\prod_{\nu=1}^{h-1}(1+v/\nu)
 \quad(0\leq r<h).
\end{equation}
Thus $e_{h,0}=1$, $e_{h,1}=H_{h-1}$, and
$e_{h,2}=\{H_{h-1}^2-H_{h-1}^{(2)}\}/2$.  In every degree the
coefficients follow from
\[
 \prod_{\nu=1}^{h-1}(1+v/\nu)
 =\exp\!\left(\sum_{k\geq1}
          \frac{(-1)^{k+1}}kH_{h-1}^{(k)}v^k\right).
\]

We use the descending-index convention
\begin{equation}\label{mx:mpl}
 \mxLi_{u_1,\ldots,u_k}(z_1,\ldots,z_k)
 =\sum_{n_1>\cdots>n_k\geq1}
       \frac{z_1^{n_1}\cdots z_k^{n_k}}
            {n_1^{u_1}\cdots n_k^{u_k}}.
\end{equation}
The first order will remain variable.  Intermediate trailing orders
may be nonpositive integers; Lemma~\ref{mx:eliminate} removes all such
orders.  For root choices $\mathbf a=(a_j)_{j\in I_+}$ with
$a_j^{p_j}=X_j$, form the words
\[
 0^{s_j-1}a_j\qquad(j\in I_+).
\]
Their shuffle, with each original position retained as a labeled
position before equal words are collected, is a finite formal sum
$\sum_w\mu_w w$.  Every resulting word has exactly $q$ nonzero
letters and ends in a nonzero letter; hence it has a unique form
\begin{equation}\label{mx:word}
 w=0^{k_1-1}b_1\,0^{k_2-1}b_2\cdots0^{k_q-1}b_q,
 \qquad k_i\geq1.
\end{equation}
The sequence $(b_1,\ldots,b_q)$ lists the selected root colors in the
order in which they occur in that shuffle.  Repeated colors do not
alter this rule or suppress their shuffle multiplicities.

\subsection{The mixed-sign reduction theorem}

\begin{theorem}[Finite mixed-sign reduction and every normal jet]
\label{mx:main}
Let the preceding data be fixed, with $q\geq1$ and $I_-\ne\varnothing$.
Put $K=\prod_{j\in I_+}p_j^{s_j-1}$.  Then
\begin{align}\label{mx:mainformula}
 Z(C)={}&K\sum_{\substack{\mathbf a\\a_j^{p_j}=X_j}}
       \sum_w\mu_w\sum_\alpha\sum_{h=1}^{M_\alpha}
           c_{\alpha,h}
       \sum_{r=0}^{h-1}e_{h,r}\sum_{\ell=0}^r
           (-1)^\ell\binom r\ell\notag\\[-1mm]
 &\hspace{10mm}\cdot
 \mxLi_{C-r+\ell,\,k_1-\ell,\,k_2,\ldots,k_q}
 \left(\alpha,\frac{b_1}\alpha,
                  \frac{b_2}{b_1},\ldots,\frac{b_q}{b_{q-1}}\right).
\end{align}
This is an identity of meromorphic functions of $C$ on the whole
complex plane.  It reduces, by finite algebra, to a sum of terms
\begin{equation}\label{mx:positiveoutput}
 A\,\mxLi_{C+v,\,t_2,\ldots,t_k}(z_1,\ldots,z_k),
 \qquad v\in\mathbb Z,\quad t_i\in\mathbb Z_{>0},\quad k\leq q+1.
\end{equation}
All coefficients and colors are independent of $C$.  Consequently
every derivative $\partial_C^N Z$ is obtained by differentiating only
the first polylogarithm order in either finite formula.  If no $X_j$
equals one, the continuation is entire and these derivatives are
ordinary entire functions.  When poles occur, the identities and
derivatives are meromorphic identities with the pole prescription in
Theorem~\ref{mx:poles} below.
\end{theorem}

\begin{proof}
The root filter, valid first as an identity of power series, is
\begin{equation}\label{mx:rootfilter}
 \mxLi_s(Xt^p)=p^{s-1}\sum_{a^p=X}\mxLi_s(at).
\end{equation}
Indeed the sum of the $n$th powers of those roots vanishes unless
$p\mid n$, and equals $pX^{n/p}$ otherwise.  Thus no choice of a
particular $p$th-root branch enters the result.

For completeness, take the iterated-integral forms
$\omega_0=dt/t$ and $\omega_a=a\,dt/(1-at)$, and let
$\mathcal H_w$ denote the corresponding integral from zero with
outermost letter first.  All words used here end in a nonzero letter,
so these integrals require no tangential-base-point regularization.
Expansion of the kernels into geometric series proves
\begin{equation}\label{mx:hyperlog}
 \mathcal H_w(t)=
 \mxLi_{k_1,\ldots,k_q}
       (b_1t,b_2/b_1,\ldots,b_q/b_{q-1}).
\end{equation}
Partitioning the product of integration simplices by all total
orderings of their integration variables proves the shuffle product,
including its multiplicities when letters coincide.  Applying this
to \eqref{mx:rootfilter} gives
\[
 P(t)=K\sum_{\mathbf a}\sum_w\mu_w\mathcal H_w(t).
\]

It remains to compute one rational-factor convolution.  Write
$H(t)=\mathcal H_w(t)=\sum_{m\geq1}h_mt^m$.  For $n\geq m$,
\begin{align}\label{mx:binomial}
 [t^{n-m}](1-\alpha t)^{-h}
 &=\alpha^{n-m}\binom{n-m+h-1}{h-1}\notag\\
 &=\alpha^{n-m}\sum_{r=0}^{h-1}e_{h,r}
                   \sum_{\ell=0}^r(-1)^\ell\binom r\ell
                            n^{r-\ell}m^\ell.
\end{align}
The diagonal $n=m$ contributes $h_n$ exactly once.  The strict part
$n>m$, followed by the nested indices already in $h_m$, is precisely
the multiple polylogarithm in \eqref{mx:mainformula}.  The diagonal
contributions from every rational partial fraction and from
$c_\infty H$ have combined coefficient
$c_\infty+\sum_{\alpha,h}c_{\alpha,h}=R(0)=0$.  Thus they cancel
exactly, leaving the displayed formula.

This reasoning proves an identity of coefficient sequences.  All
series converge locally normally in $C$ if $|X_j|<1$: the cumulative
colors in every term of \eqref{mx:mainformula} are
$(\alpha,b_1,\ldots,b_q)$, all strictly inside the disk, so the
coefficients have exponential decay times a polynomial.  For boundary
colors the same coefficient identity can first be summed in a
sufficiently far right half-plane.  Lemma~\ref{mx:eliminate} produces
\eqref{mx:positiveoutput}; each elimination decreases depth by one.
The initial depth is $q+1$, proving the asserted depth bound.
The Mellin argument in Theorem~\ref{mx:poles} proves the continuation
and the entire assertion.  Equality on the original half-plane then
implies equality meromorphically everywhere and permits every fixed
$C$ derivative.
\end{proof}

\begin{lemma}[Eliminating a nonpositive trailing order]
\label{mx:eliminate}
Any multiple polylogarithm whose first order is $C+v$ with
$v\in\mathbb Z$ and whose remaining orders are integers reduces
finitely to terms of the form \eqref{mx:positiveoutput}, with depth
no greater than its initial depth.  Coincident cumulative colors are
handled directly, without division by their difference.
\end{lemma}

\begin{proof}
Choose any nonpositive trailing order $-M$ at position $j\geq2$.
For fixed neighboring indices put
$u=n_{j-1}$ and $l=n_{j+1}$, with $n_{k+1}=0$ when $j=k$ is last.
The sum to eliminate is
\[
 \sum_{l<n<u}\beta^n n^M,\qquad \beta=z_j.
\]
For $\beta\ne1$ define the rational functions and polynomial
\begin{equation}\label{mx:Q}
 S_r(\beta)=\left(\beta\frac d{d\beta}\right)^r
                         \frac1{1-\beta},\qquad
 Q_M(x;\beta)=\sum_{r=0}^M\binom Mr x^{M-r}S_r(\beta).
\end{equation}
The geometric-series identity, or its polynomial continuation in
$\beta$, gives
\begin{equation}\label{mx:geomelim}
 \sum_{l<n<u}\beta^n n^M
   =\beta^{l+1}Q_M(l+1;\beta)-\beta^uQ_M(u;\beta).
\end{equation}
For example, it follows by subtracting the two tails
$\sum_{n\geq v}\beta^n n^M=\beta^vQ_M(v;\beta)$ when
$|\beta|<1$, and then using equality of rational functions.  At
$\beta=1$ the formula used by the algorithm is instead
\begin{equation}\label{mx:faulhaber}
 \sum_{l<n<u}n^M
   =\frac{B_{M+1}(u)-B_{M+1}(l+1)}{M+1},
\end{equation}
where $B_1(x)=x-1/2$.  Both identities hold even when $u=l+1$;
their two boundary terms then cancel.

Expand the boundary polynomials in $u$ or $l$.  A power $u^a$ lowers
the preceding order by $a$, and the factor $\beta^u$ multiplies the
preceding color by $\beta$.  Similarly, a power $l^a$ lowers the next
order and $\beta^l$ multiplies its color by $\beta$.  At the last
position the lower boundary is a scalar.  After removing position
$j$, every term is a polylogarithm of depth one less, with first order
still of the form $C+$ an integer.  Some adjacent trailing orders
may now be nonpositive; repeat the same step.  The depth strictly
decreases at each step, so the recursion terminates after at most
$k-1$ eliminations along any branch.  Every trailing order in the
final terms is positive.  Equation~\eqref{mx:faulhaber} is the exact
coincident-color branch; it never requires taking numerically unstable
limits of \eqref{mx:geomelim}.
\end{proof}

\begin{remark}[What the reduction preserves]
Only the first polylogarithm order varies with $C$.  The theorem does
not authorize differentiating the finite expression with respect to
the discrete inner indices $s_j$.  Boundary elimination replaces a
pair of adjacent colors by their product; hence its cumulative colors
are obtained by deleting entries from
$(\alpha,b_1,\ldots,b_q)$.  In particular, if every original $X_j$
is nontrivial, it cannot introduce a cumulative color equal to one.
This observation also explains why equal neighboring colors are
harmless while a genuinely uncolored original factor can change the
outer spectral poles.
\end{remark}

The two unmixed sectors fit the same proof.  If $I_-=\varnothing$,
the shuffle gives a sum of depth-$d$ polylogarithms with first order
$C+k_1$.  If $I_+=\varnothing$, the coefficient identity
\eqref{mx:binomial} gives the preceding depth-one formula
\[
 Z(C)=\sum_{\alpha,h}c_{\alpha,h}
                    \sum_{r=0}^{h-1}e_{h,r}\mxLi_{C-r}(\alpha).
\]
Thus every pattern of integer inner orders has a finite reduction.

\subsection{Continuation, poles, and normal differentiation}

Positive uncolored factors introduce logarithms.  Therefore the
finite simple-pole list from the all-nonpositive rational sector
cannot simply be copied to the mixed sector.  The following local
rule gives all the poles and their principal coefficients.

\begin{theorem}[Complete principal-part prescription]
\label{mx:poles}
For the integer-order sum \eqref{mx:Z}, including the unmixed sectors,
put
\[
 F(x)=\prod_j\mxLi_{s_j}(X_je^{-p_jx}),\qquad
 M=\sum_{\substack{s_j\leq0\\X_j=1}}(1-s_j),\qquad
 \ell_+=\#\{j:s_j>0,\ X_j=1\}.
\]
Near $x=0^+$ there is a convergent logarithmic Laurent expansion
\begin{equation}\label{mx:locallaurent}
 F(x)=\sum_{v\geq-M}\sum_{b=0}^{\ell_+} a_{v,b}x^v(\log x)^b.
\end{equation}
The meromorphic continuation of $Z$ has possible poles only at
integers $m\leq M$.  Its complete principal part at $C=m$, with
$\varepsilon=C-m$, is
\begin{equation}\label{mx:principal}
 \operatorname{PP}_{C=m}Z(C)
 =\operatorname{PP}_{\varepsilon=0}\left\{
  \frac1{\Gamma(m+\varepsilon)}
    \sum_{b=0}^{\ell_+}\frac{(-1)^bb!\,a_{-m,b}}
                        {\varepsilon^{b+1}}\right\}.
\end{equation}
Consequently the pole order is at most $\ell_++1$ for positive $m$, and
at most $\ell_+$ for nonpositive $m$.  Missing coefficients and canceled
principal parts mean removable singularities.  If no $X_j=1$,
the function is entire.  Every normal derivative is obtained by
differentiating the meromorphic function and its displayed principal
parts.
\end{theorem}

\begin{proof}
For a nontrivial color the factor is analytic at zero, with expansion
\begin{equation}\label{mx:nontrivexp}
 \mxLi_s(Xe^{-px})
   =\sum_{k\geq0}\mxLi_{s-k}(X)\frac{(-px)^k}{k!}.
\end{equation}
For $a\geq0$ the uncolored nonpositive factor has expansion
\begin{equation}\label{mx:negativeexp}
 \mxLi_{-a}(e^{-px})
   =\frac{a!}{(px)^{a+1}}
      +\sum_{k\geq0}\zeta(-a-k)\frac{(-px)^k}{k!}.
\end{equation}
For $s\geq1$ the uncolored positive factor has the integer-order
limit of Jonqui\`ere's expansion~\cite[Eq.~25.12.12]{DLMFpolylog}:
\begin{equation}\label{mx:positiveexp}
 \mxLi_s(e^{-px})
  =\frac{(-px)^{s-1}}{(s-1)!}
       \{H_{s-1}-\log p-\log x\}
    +\sum_{\substack{k\geq0\\k\ne s-1}}
       \zeta(s-k)\frac{(-px)^k}{k!}.
\end{equation}
These convergent local expansions prove \eqref{mx:locallaurent} and
provide every needed coefficient by finite multiplication.

In the initial right half-plane, termwise integration gives
\begin{equation}\label{mx:mellin}
 Z(C)=\frac1{\Gamma(C)}\int_0^\infty x^{C-1}F(x)\,dx.
\end{equation}
The integral over any $[\delta,\infty)$ is entire before multiplication
by $1/\Gamma(C)$, since $F$ decreases exponentially at infinity.
Subtract successively more terms of \eqref{mx:locallaurent} near zero.
The remainder then defines a holomorphic function on successively
larger left half-planes, and the subtracted monomials contribute
\[
 \int_0^1x^{C+v-1}(\log x)^b\,dx
       =\frac{(-1)^bb!}{(C+v)^{b+1}}.
\]
Only $v=-m$ can contribute to the principal part at $C=m$.
The reciprocal Gamma function is nonzero at positive integers and
has a simple zero at each nonpositive integer.  This proves
\eqref{mx:principal} and the pole-order assertions.  If there are no
uncolored factors, then $M=\ell_+=0$; all potential Mellin poles cancel
against those Gamma zeros.  Local normal convergence of the subtracted
integrals on compact sets also justifies every fixed normal
derivative.

To verify continuation of each nested term separately, its coefficient
kernel, with the first denominator omitted, is
\[
 \sum_{n_1>\cdots>n_k\geq1}
 \frac{(z_1t)^{n_1}z_2^{n_2}\cdots z_k^{n_k}}
      {n_2^{t_2}\cdots n_k^{t_k}}
 =\frac{z_1t}{1-z_1t}
     \mxLi_{t_2,\ldots,t_k}(z_1z_2t,z_3,\ldots,z_k).
\]
For depth one the kernel is just $z_1t/(1-z_1t)$.  The positive
trailing indices make the right side a rational function times an
iterated integral.  Induction in its word length gives a convergent
finite-logarithmic local expansion at $t=1$.  Applying the same Mellin
subtractions with parameter $C+v$ proves its meromorphic continuation.
If none of its cumulative colors is one, its kernel is analytic at
$t=1$, and the continuation is entire.  As noted after
Lemma~\ref{mx:eliminate}, all these cumulative colors come from the
original root sets.
\end{proof}

Formula~\eqref{mx:principal} is fully finite at every specified pole.
For clarity, the Gamma factors needed there can be generated by
\begin{align}\label{mx:gammajets}
 \frac1{\Gamma(m+\varepsilon)}
 &=\frac1{(m-1)!}\exp\left\{
   (\gamma-H_{m-1})\varepsilon+
   \sum_{k\geq2}\frac{(-1)^{k+1}}k
       (\zeta(k)-H_{m-1}^{(k)})\varepsilon^k\right\},\quad m\geq1,
 \notag\\
 \frac1{\Gamma(-N+\varepsilon)}
 &=(-1)^NN!\varepsilon\exp\left\{
   (\gamma-H_N)\varepsilon+
   \sum_{k\geq2}\frac{(-1)^{k+1}\zeta(k)-H_N^{(k)}}k
                                        \varepsilon^k\right\}.
\end{align}
Only finitely many coefficients are required.  This controls all
principal coefficients through ordinary zeta values, generalized
harmonic numbers, integer-order polylogarithms, and logarithms of
the weights.  It does not assert that the holomorphic Laurent
coefficients reduce to depth one.

Inside the unit polydisk the functions are already entire in $C$.
When all boundary colors are nontrivial, their continuation agrees
locally uniformly in $C$ with the total-degree Abel limit
$X_j\mapsto\rho^{p_j}X_j$, $\rho\uparrow1$.  This follows from
the same Taylor subtractions, uniformly near those fixed colors.
If uncolored factors occur, an Abel limit outside the convergence
half-plane need not exist.  In that case the prescribed value is the
meromorphic continuation of the fixed-color sum, and a pole requires
a Laurent coefficient or other explicitly stated restriction.

\subsection{Explicit mixed and weighted identities}

For compactness write
\[
 Z_{s_1,\ldots,s_d}(C;z)
  =\sum_{m_1,\ldots,m_d\geq1}
      \frac{z^{m_1+\cdots+m_d}}
       {m_1^{s_1}\cdots m_d^{s_d}(m_1+\cdots+m_d)^C}.
\]
The following identities hold first for $|z|<1$, and thereafter with
the fixed-color continuation just specified.  From
$\mxLi_1(zt)^2=2\mxLi_{1,1}(zt,1)$ and
$\mxLi_0(zt)=-1+(1-zt)^{-1}$, one obtains
\begin{equation}\label{mx:110}
 Z_{1,1,0}(C;z)=2\mxLi_{C,1,1}(z,1,1).
\end{equation}
The repeated pole
\[
 \mxLi_{-1}(zt)=(1-zt)^{-2}-(1-zt)^{-1}
\]
gives, after the coincident-color elimination,
\begin{align}\label{mx:11minus1}
 Z_{1,1,-1}(C;z)={}&2\mxLi_{C-1,1,1}(z,1,1)
                   -2\mxLi_{C-1,1}(z,1)
                   +2\mxLi_{C,1}(z,1)\notag\\
                 &+2\mxLi_{C-1}(z)-2\mxLi_C(z).
\end{align}
Indeed the intermediate term is
$2\{\mxLi_{C-1,1,1}-\mxLi_{C,0,1}\}$, and direct summation of
the middle index proves
\[
 \mxLi_{a,0,1}(z,1,1)
 =\mxLi_{a-1,1}(z,1)-\mxLi_{a,1}(z,1)
      -\mxLi_{a-1}(z)+\mxLi_a(z).
\]
Every normal derivative of \eqref{mx:11minus1}, including at $C=0$
when $z\ne1$, follows by replacing each polylogarithm with its
corresponding first-order derivative.  This gives the whole normal
jet, rather than an isolated negative-integer value identity.

A four-variable example has depth at most three, as predicted:
\begin{align}\label{mx:1100}
 Z_{1,1,0,0}(C;z)={}&2\mxLi_{C-1,1,1}(z,1,1)
                    -2\mxLi_{C,1,1}(z,1,1)\notag\\
                  &-2\mxLi_{C-1,1}(z,1)
                    +2\mxLi_{C,1}(z,1)
                    +2\mxLi_{C-1}(z)-2\mxLi_C(z).
\end{align}
Here the rational factor is
$(zt)^2/(1-zt)^2=1-2/(1-zt)+(1-zt)^{-2}$.

An unequal-weight example, with both root filtering and a color
coincidence, is
\begin{equation}\label{mx:weightedexample}
 W(C)=\sum_{m,n,k\geq1}
       \frac{4^{-m}3^{-n}9^{-k}}{m^2n(2m+n+2k)^C}.
\end{equation}
It is entire in $C$.  Since
\[
 \mxLi_2(t^2/4)=2\{\mxLi_2(t/2)+\mxLi_2(-t/2)\},\qquad
 \frac{t^2/9}{1-t^2/9}
   =-1+\frac1{2(1-t/3)}+\frac1{2(1+t/3)},
\]
the three shuffles of the words $0a$ and $1/3$ give
\begin{align}\label{mx:weightedanswer}
 W(C)=\sum_{\substack{a\in\{1/2,-1/2\}\\
                     \delta\in\{1/3,-1/3\}}}
 \biggl\{&\mxLi_{C,1,2}(\delta,1/(3\delta),3a)
       +\mxLi_{C,2,1}(\delta,1/(3\delta),3a)\notag\\
       &+\mxLi_{C,2,1}(\delta,a/\delta,1/(3a))\biggr\}.
\end{align}
All cumulative colors in these twelve terms have modulus less than
one, so no continuation convention is needed to differentiate this
identity at any $C$.

The uncolored case of \eqref{mx:110} illustrates why the pole theorem
is necessary.  Near zero its Mellin kernel is
\[
 \mxLi_0(e^{-x})\mxLi_1(e^{-x})^2
 =x^{-1}\log^2x-\tfrac12\log^2x-\log x
   +x\{\tfrac1{12}\log^2x+\tfrac7{12}\log x+\tfrac14\}
   +O(x^2(1+|\log x|^2)).
\]
Consequently
\begin{align}\label{mx:examplepoles}
 \operatorname{PP}_{C=1} Z_{1,1,0}(C;1)
  &=\frac2{(C-1)^3}+\frac{2\gamma}{(C-1)^2}
                         +\frac{\gamma^2-\zeta(2)}{C-1},\notag\\
 \operatorname{PP}_{C=0} Z_{1,1,0}(C;1)
  &=-\frac1{C^2}+\frac{1-\gamma}{C},\notag\\
 \operatorname{PP}_{C=-1} Z_{1,1,0}(C;1)
  &=-\frac1{6(C+1)^2}+\frac{3/4-\gamma/6}{C+1}.
\end{align}
In fact there are infinitely many nonpositive poles: the coefficient
of $\log^2x$ in this kernel is exactly
$\mxLi_0(e^{-x})$, whose positive odd powers have nonzero Bernoulli
coefficients.  This observation concerns the mixed extension; it is
not a correction of the earlier theorem restricted to rational
generating functions.

\subsection{A compatible primitive hierarchy and exact checks}

With $r_n$ from \eqref{mx:factor}, introduce an additional variable
$z$ by $Z(C;z)=\sum_{n\geq1}r_nz^nn^{-C}$.  Equivalently this
replaces $X_j$ by $X_jz^{p_j}$.  In every term of the finite reduction
only the first polylogarithm color acquires the factor $z$.  Thus
\begin{equation}\label{mx:primitives}
 \int_0^z\partial_C^N Z(C;t)\frac{dt}{t}
        =\partial_C^N Z(C+1;z),\qquad
 z\frac d{dz}\partial_C^N Z(C;z)
        =\partial_C^N Z(C-1;z).
\end{equation}
The primitive is fixed by its zero value at $z=0$.  These identities
follow by local normal convergence in the disk and continue along
any common path avoiding the coefficient-function singularities.

The accompanying program performs 624 exact rational checks with
SymPy; its source is \texttt{verify\_mixed\_reduction.py}.
There are 484 coefficient
checks: for eleven independent fixtures, each of the first 22
coefficients is computed directly from the original product and
compared both with \eqref{mx:mainformula} and with the fully eliminated
positive-trailing-order expression.  The fixtures include four
original variables, three positive factors, unequal weights, repeated
rational poles, coincident and opposite unit colors, and nonpositive
intermediate trailing indices.  The remaining 140 checks verify the
finite boundary summations for degrees zero through six, including
empty intervals and $\beta=1$.  No floating-point tolerance is used.
The coefficient checks establish finite algebraic diagnostics;
meromorphic continuation and differentiation are supplied by the
proofs, not inferred from these tests.  The exact outputs and fixture
descriptions are recorded in \texttt{mixed\_reduction\_checks.json}.

The remaining question is arithmetic simplification of the resulting
positive-depth normal jets.  The depth bound is a guaranteed upper
bound, not an independence statement or a proof of minimality.
At roots of unity, one can now ask which symmetries or distribution
projections cancel these residual jets and leave only depth-one
Hurwitz-zeta, Gamma, and generalized-Stieltjes data.  Such a
cancellation problem is distinct from the finite mixed-sign closure
proved here.

\section{Verification, source corrections, and further research}
\label{sec:audit}

\subsection{What the checks establish}

The mathematical proofs above are independent of floating-point
agreement. The supplied programs exercise three different layers of
the argument.

\paragraph{Exact finite algebra.}
The cubic symbolic script treats $\gamma,L,\zeta(2),\zeta(3)$,
the needed zeta derivatives, and $\Omega_3,J_3$ as formal symbols.
It reconstructs the four quadratic coefficients of $H$ from the
four stated restrictions and reduces the residual in
\eqref{ray:main} to the zero rational function. Its curve checks
compare direct truncated quotient expansion with
\eqref{curve:differentials}; they retain all path coefficients.
The mixed-order script independently constructs the original
coefficient convolution, the shuffle/partial-fraction formula, and
the output after eliminating nonpositive trailing orders. Its 624
exact rational checks include coincident colors and repeated poles.
A separate local-coefficient script checks the displayed Mellin
principal parts and Gamma jets and the finite harmonic
summation-by-parts algebra in Section~\ref{bridge:sec}. In total,
the three exact suites record 760 passed checks: 16 cubic/ray
identities, 624 mixed-order checks, and 120 local-coefficient checks.

\paragraph{A certified scalar evaluation.}
The code \texttt{verify\_cubic\_harmonic\_bridge.py} combines
Arb ball arithmetic with \eqref{bridge:tail-bound}, using
$N=50$, $K=24$ and 100 decimal digits of working precision.
The stored enclosure is
\[
 Q\in 1.9662627343915343406643753384116314238258689764916868267443076
       \ \pm\ 4.86\cdot10^{-62}.
\]
The analytic tail contribution alone is below $1.26\cdot10^{-62}$.
The reported enclosure includes rounding of its displayed center.
The corresponding enclosures for $\eta(1)$ and $\Omega_3$ are
obtained using the proved exact identities. They are not independent
proofs of those identities.

\paragraph{Independent numerical diagnostics.}
Ordinary high-precision quadrature of the regularized digamma cube
and the harmonic Mellin integral agree through approximately 70
decimal places. A different evaluator computes $\T$ by integrating
Jonqui\`ere's series on $(0,1)$ and an incomplete-Gamma tail on
$(1,\infty)$. Symmetric differences with polynomial extrapolation
check the third derivatives on six additional rays, including
$(1,-2,3)$. Their largest recorded absolute discrepancy is below
$5.7\cdot10^{-37}$. An independent computation of $I_2$ gives the
same second cubic coordinate through more than 30 decimal places.
These latter calculations use non-interval arithmetic and truncated
series; their digits are diagnostics, not certified error estimates.
The split at one is strictly inside the Jonqui\`ere radius $2\pi$.

The package records exact program versions, numerical parameters,
machine-readable outputs, and build instructions. No finite test
suite certifies the analytic convergence arguments: those arguments
are supplied in the text. The exact symbolic checks certify the
specified finite algebra, rather than all possible extensions of a
theorem.

\subsection{A source-status ledger for integration}

The source snapshot is pinned by its full commit identifier, with
blob hashes in \path{provenance/repository_inventory.json}.
All 80 manuscript chapter files and all 11 incoming archives matched
their recorded Git blob hashes. This establishes the identity of the
source corpus; it is not a claim that every proof in that corpus has
received a new exhaustive review. The mathematical audit concentrated
on the harmonic normalization, the cubic completion, the local
Tornheim decomposition, and the rational/colored reduction machinery
used by the present results.

\begin{longtable}{@{}>{\raggedright\arraybackslash}p{0.25\textwidth}>{\raggedright\arraybackslash}p{0.46\textwidth}>{\raggedright\arraybackslash}p{0.23\textwidth}@{}}
\toprule
Existing target & Result in this report & Status to record\\
\midrule
\endhead
Coincident report, Q1: third diagonal derivative
& Theorem~\ref{thm:diag-harmonic} and the absolutely convergent
harmonic series \eqref{bridge:arithmetic}.
& Reduced to $\eta(1)$; finite ordinary-constant closure remains open.\\[5pt]
Coincident report, Q2: asymmetric cubic directions
& Theorem~\ref{thm:cubic-rays}, with two fixed coordinates and
explicit coefficients; $\kappa$ also has the convergent integral
\eqref{ray:I2}.
& Complete spanning formula proved; arithmetic minimality unproved.\\[5pt]
Coincident report, Q11: mixed integer inner signs
& Theorem~\ref{mx:main}, including arbitrary integer weights,
coincident colors, repeated poles and all outer-order derivatives.
& Finite closure proved, with a guaranteed depth bound.\\[5pt]
Strict/inclusive harmonic conventions
& \eqref{eq:strict-inclusive} is retained exactly, including
the sign and the $\zeta'(2,a)$ shift.
& Existing normalization confirmed.\\[5pt]
Gaussian $S_6$ and revised $S_8$ candidates
& No exact target-vector certificate is supplied here.
& Remain open; no theorem label should be promoted.\\
\bottomrule
\end{longtable}

\subsection{Corrections and qualifications worth preserving}

No new false theorem was found in the specific repository results
used here. Several distinctions are nevertheless necessary for a
correct integration.

\begin{enumerate}
\item \textbf{Keep the full resonant numerators.}
The cubic completion \eqref{diag:completed} subtracts
$3h(t)^2/t$ and adds $3(1-t)\zeta(2-t)/(2t)$.
Replacing their numerators by their constant terms changes the
answer. This is an existing warning in the coincident report,
rederived here to protect the new harmonic bridge from a finite
normalization error.
\item \textbf{Use the mixed pole rule only in its stated domain.}
The rational sector has a finite simple-pole description. Its
extension with positive uncolored inner orders can have higher
pole order and infinitely many nonpositive poles. The explicit
$Z_{1,1,0}(C;1)$ example proves this difference. The earlier
rational theorem is not wrong; extending its pole statement without
the logarithmic correction would be wrong.
\item \textbf{Separate coordinate rank from arithmetic independence.}
Two formal quadratic coordinates of $H$ survive the exact slice
constraints, with slope coefficients $abc$ and $c\Delta$.
The theorem determines that reconstruction problem. It does not
show that $\eta(1)$ or $\kappa$ lies outside a chosen field of
ordinary constants, or that they are independent of each other.
\item \textbf{Restrict before differentiating.}
An admissible Tornheim ray is holomorphic at zero, while the full
three-variable function is meromorphic there. The values $1/3$ on
the diagonal and $5/12$ on the $C$ axis are consistent. A curve
tangent to $A+C=0$ or $B+C=0$ is outside
Theorem~\ref{thm:curve}; its value is not assigned by cancellation
of a zero denominator in the displayed ray formula.
\end{enumerate}

The incoming report already records corrections to the inspected
Borwein--Dilcher author PDF \cite{IncomingCoincident,BorweinDilcher}.
They should retain that attribution. In the integer Jonqui\`ere
formula, the convention is $H_0=0$, and the coefficient of
$x^{p-1}\log x$ is
\[
 B_p=\frac{(-1)^p}{\Gamma(p)},\qquad p\ge1,
\]
without an additional factor $H_{p-1}$. This follows at once from
\eqref{mx:positiveexp}; $p=3$ gives $-1/2$.
For the ordinary convergent local power series, the splitting point
must satisfy $0<\theta<2\pi$. These are previously documented
corrections to the inspected author version, not newly discovered
errors or assertions about every published version.

\subsection{Proposed further research questions}

The following problems are specific next steps. They are questions,
not conjectured equalities unsupported by calculation.

\paragraph{1. Determine the arithmetic content of the harmonic bridge.}
Fix a constant algebra before attempting reduction, for example one
generated over $\mathbb Q$ by ordinary Stieltjes constants, ordinary
zeta derivatives at integers, logarithms and specified Barnes values.
Does $\eta(1)$ have a finite expression in that algebra? The bridge
gives three analytic realizations of the same coordinate: an ordered
double-zeta Laurent coefficient, a convergent logarithmic harmonic
series, and a normalized digamma-cube integral. Seek a functional
identity relating two realizations without reintroducing the same
coefficient under another name.

\paragraph{2. Relate the second cubic coordinate to harmonic Laurent data.}
The constant $\kappa$ is a fixed asymmetric double-zeta derivative
and also the convergent moment \eqref{ray:I2} plus classical terms.
Can an ordered Hurwitz functional equation carry that derivative
from the origin to a harmonic Stieltjes coefficient at $(1,1)$?
A successful identity should state the exact finite correction and
the direction of continuation. The question of a relation between
$\kappa$ and $\eta(1)$ is distinct from the proved two-coordinate
spanning formula.

\paragraph{3. Classify useful cancellations among cubic directions.}
The conic $\Delta=0$ removes $\kappa$ from a single ray. More
generally, a finite combination $\sum_j r_jW_{v_j}'''(0)$ has no
displayed nonclassical coordinate whenever
$\sum_jr_ja_jb_jc_j=0$ and
$\sum_jr_jc_j\Delta_j=0$. Classify short rational configurations
with additional permutation, distribution, or character symmetry.
This is a finite exact linear problem and can produce concrete new
single-zeta identities without an independence assumption.

\paragraph{4. Construct the complete quartic directional layer.}
At fourth order the cubic homogeneous part of $H$ enters. Determine
its exact dimension after symmetry and the axis/stuffle restrictions,
then choose independent anchor slices and derive the analogue of
\eqref{ray:main}. The first goal is a proved finite coordinate
formula with convergent integral representatives. A finite
ordinary-zeta evaluation should be a separate theorem if achieved.

\paragraph{5. Extend the cubic bridge to a common Hurwitz shift.}
For $a>0$, the coefficient $\eta(a)$ has exact shift and derivative
relations in the ordered report. Apply the partial-fraction method
to shifted poles and compare with the primitive
$\mathcal P_3(a+1)-\mathcal P_3(a)$. Determine the explicit
boundary terms and whether products of shifted digamma and
polygamma functions give a useful identity involving $\eta(a)$.
The unit-coordinate finite part at a singular endpoint and an
ordinary integral over a positive shifted interval must stay distinct.

\paragraph{6. Obtain a general power-to-harmonic formula.}
For each $k\ge4$, the principal parts of $\psi(z)^k$ are finite
generalized-harmonic polynomials. The two-point contour argument
still suggests a normally convergent subtraction. Carry out the
finite summation-by-parts reduction explicitly for $k=4$ and identify
the required harmonic Laurent coefficients or higher-depth data.
The target is a compact identity with every constant fixed, rather
than an unspecified Mittag--Leffler entire remainder.

\paragraph{7. Find depth-reducing projections in the mixed sector.}
At root-of-unity colors, the finite reduction outputs a finite list
of nested functions and their first-order derivatives. Construct
exact matrices for character projections and distribution identities
on that list. Their kernels can identify combinations whose depth
drops to one, where Hurwitz-zeta, Gamma and generalized-Stieltjes
formulas become available. Any claimed minimal depth must specify
the relation space being used.

\paragraph{8. Evaluate selected holomorphic Laurent coefficients.}
Theorem~\ref{mx:poles} gives all principal parts but deliberately
does not promise finite depth-one formulas for the regular terms.
Begin with $Z_{1,1,0}(C;1)$ at $C=1,0,-1$: derive explicit
convergent kernels for its constant and linear Laurent coefficients
and relate them to depth-three harmonic Stieltjes conventions.
This separates the already solved local pole algebra from new
global arithmetic.

\paragraph{9. Normalize ordinary primitives under reflection and multiplication.}
The primitive $\mathcal P_3$ fixes its integration constant at the
singular origin, while \eqref{mx:primitives} fixes a holomorphic
Euler primitive at zero. Determine exact reflection and multiplication
relations for these normalizations, including their polynomial and
logarithmic correction terms. At low orders, compare with Barnes
multiple-Gamma primitives only after matching the additive constants
and branches.

\paragraph{10. Resolve tangential spectral paths.}
The transverse formula requires $(a+c)(b+c)\ne0$. For a curve
whose first nonzero contact with one polar divisor occurs at order
$r>1$, determine when the numerator cancels enough terms to leave
a removable restriction, and which path jets fix the resulting
finite value. This is a divisibility and coefficient-extraction
problem, not an instruction to apply the transverse formula at a
zero denominator.

\paragraph{11. Return to the exact Gaussian targets.}
The frozen $S_6$ and revised $S_8$ vectors still need exact
certificates in proved relation spaces. Use the current reduction
and projection machinery to propose additional rows, proving the
regularization prescription for each row. A numerical integer
relation is a candidate; a successful proof must combine exact
rows to the exact target. An unsuccessful restricted search should
report its precise space and a separating functional, not an
unrestricted impossibility claim.

\paragraph{12. Formalize the finite proof interfaces.}
The partial-fraction extraction, shuffle multiplicities, finite
antidifferences, Gamma jets, and cubic coefficient elimination are
all terminating algebra. Separate them from analytic lemmas for
normal convergence, Mellin continuation and endpoint constant
terms. A formal implementation can then check the finite identities
while explicitly declaring the analytic hypotheses on which their
applications depend.

These questions preserve the emphasis on exact identities and
normalizations. Error bounds appear in this report only to certify
the numerical evaluation of an already proved identity.

\begin{thebibliography}{99}
\raggedright
\addcontentsline{toc}{section}{References}

\bibitem{Repo}
V. Reshetnikov and the ProveIt research programme,
\emph{Polylogarithms and their Arithmetic Bridges}, manuscript and
incoming research corpus, snapshot \sourcecommit, accessed
11 October 2026 (UTC).
\href{https://github.com/VladimirReshetnikov/ProveIt/tree/5a790187c8e186e41e2b990b4941cb7a1a3c7b6b/Analysis/Polylogarithms/docs/manuscript}{Pinned manuscript directory};
\href{https://github.com/VladimirReshetnikov/ProveIt/tree/5a790187c8e186e41e2b990b4941cb7a1a3c7b6b/docs/incoming}{pinned incoming directory}.

\bibitem{IncomingCoincident}
ProveIt research continuation,
\emph{Coincident Stieltjes Products and Directional Zeta Identities:
Cubic finite parts, distributional contact terms, coordinate
invariance, and exact spectral reductions},
10 October 2026, incoming archive
\texttt{ProveIt\_Coincident\_Stieltjes\_Directional\_Zeta\_2026-10-10.zip}
at the revision in \cite{Repo}.
Relevant source files: \texttt{sections/02\_cubic.tex},
\texttt{05\_spectral.tex}, \texttt{06\_rays.tex},
\texttt{08b\_external\_audit.tex}, and \texttt{09\_research.tex}.
The question numbers in this report refer to that last file.

\bibitem{IncomingOrdered}
ProveIt research continuation,
\emph{Ordered Hurwitz Resonances and Cubic Stieltjes Products:
Exact Laurent coefficients, Gamma generating functions, sharp curve
dependence, and a polylogarithmic Mellin bridge},
10 October 2026, incoming archive
\texttt{ProveIt\_Ordered\_Hurwitz\_Resonances\_2026-10-10.zip}
at the revision in \cite{Repo}.
Relevant source files: \texttt{sections/02\_ordered.tex},
\texttt{04\_harmonic\_polylog.tex}, and \texttt{05\_cubic.tex}.
The archive source fixes the strict harmonic coefficient used here.

\bibitem{Kargin}
L. Karg\i n, A. Dil, M. Cenkci and M. Can,
\emph{On the Stieltjes constants with respect to harmonic zeta
functions}, Journal of Mathematical Analysis and Applications
\textbf{525} (2023), article 127302.
\url{https://doi.org/10.1016/j.jmaa.2023.127302}.
Author preprint: \url{https://arxiv.org/abs/2304.03517}.

\bibitem{MOW}
K. Matsumoto, T. Onozuka and I. Wakabayashi,
\emph{Laurent series expansions of multiple zeta-functions of
Euler--Zagier type at integer points}, Mathematische Zeitschrift
\textbf{295} (2020), 623--642.
\url{https://doi.org/10.1007/s00209-019-02337-2}.
Author preprint: \url{https://arxiv.org/abs/1601.05918}.

\bibitem{BorweinDilcher}
J. M. Borwein and K. Dilcher,
\emph{Derivatives and fast evaluation of the Tornheim zeta function},
The Ramanujan Journal \textbf{45} (2018), 413--432.
\url{https://doi.org/10.1007/s11139-017-9890-9}.
Inspected author version, titled \emph{Derivatives and fast evaluation
of the Witten zeta function}:
\url{https://carmamaths.org/jon/omega3.pdf}.
The version-specific corrections cited in Section~\ref{sec:audit}
were already documented in \cite{IncomingCoincident}.

\bibitem{Connon}
D. F. Connon, \emph{On an integral involving the digamma function},
preprint (2012), arXiv:1212.1432.
\url{https://arxiv.org/abs/1212.1432}.
Used for context on digamma integrals, not as a source of the
cubic--harmonic theorem proved in this report.

\bibitem{DLMFpolylog}
NIST Digital Library of Mathematical Functions,
Section 25.12, \emph{Polylogarithms}, especially Eq.~25.12.12.
\url{https://dlmf.nist.gov/25.12}. Accessed 11 October 2026.

\bibitem{DLMFhurwitz}
NIST Digital Library of Mathematical Functions,
Section 25.11, \emph{Hurwitz zeta function}.
\url{https://dlmf.nist.gov/25.11}. Accessed 11 October 2026.

\bibitem{DLMFgamma}
NIST Digital Library of Mathematical Functions,
Chapter 5, \emph{Gamma function}, especially Sections 5.7, 5.9,
5.11 and 5.15.
\url{https://dlmf.nist.gov/5}. Accessed 11 October 2026.

\bibitem{mx:Panzer}
E. Panzer, \emph{Algorithms for the symbolic integration of
hyperlogarithms with applications to Feynman integrals},
Computer Physics Communications \textbf{188} (2015), 148--166.
\url{https://doi.org/10.1016/j.cpc.2014.10.019}.
Author preprint: \url{https://arxiv.org/abs/1403.3385}.
Sections 2.1--2.3 supply classical background on shuffle products,
rational prefactors and local logarithmic expansions.

\bibitem{mx:Zhao}
J. Zhao, \emph{A note on colored Tornheim's double series},
Integers \textbf{10} (2010), A59, 879--882.
\url{https://arxiv.org/abs/0907.5106}.

\end{thebibliography}

\end{document}
